\Opensolutionfile{ans}[ans/ansTN-QuocHoc122]

\noindent \textbf{PHẦN I. Câu trắc nghiệm nhiều phương án lựa chọn. Thí sinh trả lời từ câu 1 đến câu 12. Mỗi câu hỏi thí sinh chỉ chọn một phương án}
\vspace{0.2cm}

\begin{ex}	\macau{ID: VL12-QH122-C01}Tính chất nào sau đây \textbf{không} phải là tính chất của chất ở thể khí?
	\choice
	{Có lực tương tác phân tử nhỏ hơn lực tương tác phân tử ở thể rắn và thể lỏng}
	{\True Có hình dạng và thể tích riêng xác định}
	{Có thể nén được dễ dàng}
	{Có các phân tử chuyển động hoàn toàn hỗn độn}
	\loigiai{
		Các phân tử ở thể khí ở xa nhau, lực tương tác rất yếu nên chuyển động hoàn toàn hỗn loạn không ngừng về mọi phía, dẫn đến chất khí không có hình dạng và thể tích riêng xác định mà luôn chiếm toàn bộ thể tích của bình chứa. Do đó nhận định B sai}
\end{ex}

\begin{ex}	\macau{ID: VL12-QH122-C02}Một quả bóng có khối lượng $100\text{ g}$ rơi từ độ cao $10\text{ m}$ xuống sân và nảy ngược lên đạt độ cao $7\text{ m}$. Lấy gia tốc trọng trường $g=9,8\text{ m/s}^2$. Độ biến thiên nội năng của quả bóng và sàn trong quá trình trên bằng
	\choice
	{$6,86\text{ J}$}
	{\True $2,94\text{ J}$}
	{$294\text{ J}$}
	{$3,00\text{ J}$}
	\loigiai{
		- Quy đổi khối lượng quả bóng: $m = 100\text{ g} = 0,1\text{ kg}$.\\
		- Cơ năng của quả bóng tại vị trí thả rơi ban đầu: $W_1 = mgh_1 = 0,1 \cdot 9,8 \cdot 10 = 9,8\text{ J}$.\\
		- Cơ năng của quả bóng tại vị trí nảy lên cao nhất: $W_2 = mgh_2 = 0,1 \cdot 9,8 \cdot 7 = 6,86\text{ J}$.\\
		- Theo định luật bảo toàn và chuyển hóa năng lượng, phần cơ năng vĩ mô bị hao hụt do va chạm đã chuyển hóa trọn vẹn thành nhiệt năng làm tăng nội năng của hệ bóng và sân:
		$$\Delta U = W_1 - W_2 = 9,8 - 6,86 = 2,94\text{ J}$$}
\end{ex}

\begin{ex}	\macau{ID: VL12-QH122-C03}Một vật khối lượng $m$, có nhiệt dung riêng $c$, nhiệt độ đầu và cuối lần lượt kí hiệu là $t_1$ và $t_2$. Công thức $Q=cm(t_2-t_1)$ được dùng để xác định đại lượng nào?
	\choice
	{\True nhiệt lượng}
	{nhiệt năng}
	{nội năng}
	{năng lượng cơ học}
	\loigiai{
		Biểu thức lượng năng lượng nhiệt biến đổi $Q = mc(t_2 - t_1)$ dùng để định lượng giá trị nhiệt lượng mà vật thu vào (nếu $t_2 > t_1$) hoặc tỏa ra (nếu $t_2 < t_1$) khi thay đổi nhiệt độ}
\end{ex}

\begin{ex}	\macau{ID: VL12-QH122-C04}Nhiệt lượng cần thiết để nung nóng chảy hoàn toàn khối lượng $m_0\text{ (kg)}$ vàng tại nhiệt độ nóng chảy của nó là $Q$. Nếu chỉ cung cấp $75\%$ lượng nhiệt lượng nói trên thì khối lượng vàng chưa bị nóng chảy tương ứng là $10\text{ g}$. Giá trị của $m_0$ ban đầu là
	\choice
	{$20\text{ g}$}
	{$10\text{ g}$}
	{\True $40\text{ g}$}
	{$30\text{ g}$}
	\loigiai{
		- Lượng nhiệt lượng cung cấp tỉ lệ thuận với khối lượng chất nóng chảy chuyển pha hoàn toàn.\\
		- Khi chỉ cung cấp $75\%Q$, hệ chỉ nóng chảy được $75\%m_0$, lượng khối lượng vàng còn tồn tại ở thể rắn chưa bị nóng chảy chiếm tỉ lệ:
		$$\Delta m = (100\% - 75\%) \cdot m_0 = 0,25m_0 = 10\text{ g} \Rightarrow m_0 = \frac{10}{0,25} = 40\text{ g}$$}
\end{ex}

\begin{ex}	\macau{ID: VL12-QH122-C05}Điều nào sau đây là \textbf{sai} khi thuyết minh về mô hình cấu tạo chất khí?
	\choice
	{\True Các nguyên tử hay phân tử chuyển động nhiệt càng nhanh thì nhiệt độ của vật càng thấp}
	{Các nguyên tử, phân tử luôn thực hiện chuyển động hỗn loạn không ngừng về mọi phía}
	{Các nguyên tử, phân tử đồng thời tồn tại cả lực hút phân tử và lực đẩy phân tử}
	{Các chất trong tự nhiên đều được cấu tạo từ các hạt riêng biệt gọi là nguyên tử, phân tử}
	\loigiai{
		Nhiệt độ là thước đo động năng trung bình của các phân tử. Do đó, các nguyên tử hay phân tử cấu tạo nên vật chuyển động nhiệt càng nhanh thì nhiệt độ của vật phải càng cao, nhận định A nói "càng thấp" là sai}
\end{ex}

\begin{ex}	\macau{ID: VL12-QH122-C06}Các chất rắn ở dạng trạng thái vô định hình sở hữu đặc điểm và tính chất vật lí nào sau đây?
	\choice
	{có tính dị hướng hình học}
	{\True có nhiệt độ nóng chảy không xác định cố định}
	{có cấu trúc mạng tinh thể tuần hoàn}
	{có dạng hình học xác định trong không gian}
	\loigiai{
		Chất rắn vô định hình (như thủy minh, nhựa đường, hắc ín) không có cấu trúc tinh thể tuần hoàn, mang thuộc tính đẳng hướng và đặc biệt không có nhiệt độ nóng chảy xác định, chúng bị mềm đi liên tục khi tăng nhiệt độ}
\end{ex}

\begin{ex}	\macau{ID: VL12-QH122-C07}Đưa một chiếc cốc chứa nước lạnh ra môi trường ngoài trời nắng nóng thì sau một thời gian thấy xuất hiện một lớp nước lỏng bám dính ở bề mặt ngoài thành cốc. Đó là kết quả của hiện tượng chuyển pha nào?
	\choice
	{bay hơi}
	{\True ngưng tụ}
	{nóng chảy}
	{thăng hoa}
	\loigiai{
		Hơi nước tồn tại sẵn trong không khí khí quyển xung quanh khi chuyển động va chạm tiếp xúc với thành cốc lạnh sẽ bị mất nhiệt độ đột ngột, chuyển pha trạng thái từ thể khí sang thể lỏng bám ngoài thành cốc. Đây là hiện tượng ngưng tụ}
\end{ex}

\begin{ex}	\macau{ID: VL12-QH122-C08}Hệ thức lượng của định luật Nguyên lí I nhiệt động lực học $\Delta U = A + Q$ khi xét điều kiện hằng số $Q > 0$ và $A < 0$ mô tả chính xác tiến trình biến đổi nội năng nào dưới đây của hệ?
	\choice
	{Hệ thực hiện truyền tỏa nhiệt lượng ra ngoài}
	{Không có đủ dữ kiện động học để đưa ra kết luận}
	{\True Hệ thực hiện thu nhận nhiệt lượng từ ngoài và đồng thời thực hiện sinh công lên ngoại cảnh}
	{Hệ thực hiện nhận công cơ học từ ngoại lực}
	\loigiai{
		Theo quy ước dấu của Nguyên lí I nhiệt động lực học:
		- Giá trị nhiệt lượng $Q > 0$ chứng tỏ hệ đang thu nhận nhiệt lượng từ môi trường.\\
		- Giá trị công đại số $A < 0$ chứng tỏ hệ đang thực hiện sinh công (lực đẩy thực hiện công) hướng ra ngoại cảnh}
\end{ex}

\begin{ex}	\macau{ID: VL12-QH122-C09}Một viên đạn đại bác có khối lượng $10\text{ kg}$ khi bay rơi tới vạch đích chạm sát đất đạt giá trị vận tốc bằng $54\text{ km/h solos}$. Nếu giả định toàn bộ phần động năng vĩ mô của nó chuyển hóa trọn vẹn thành nội năng tiêu hao tại chỗ, thì lượng nhiệt lượng tỏa ra lúc va chạm vào khoảng
	\choice
	{$2250\text{ J}$}
	{\True $1125\text{ J}$}
	{$14580\text{ J}$}
	{$7290\text{ J}$}
	\loigiai{
		- Đổi đơn vị vận tốc tức thời của viên đạn sang hệ mét trên giây chuẩn SI: $v = 54\text{ km/h} = \frac{54}{3,6} = 15\text{ m/s}$.\\
		- Động năng chuyển động vĩ mô của viên đạn đại bác ngay trước va chạm:
		$$W_{\text{đ}} = \frac{1}{2}m v^2 = \frac{1}{2} \cdot 10 \cdot 15^2 = 5 \cdot 225 = 1125\text{ J}$$
		- Vì toàn bộ động năng chuyển hóa thành nội năng, lượng nhiệt lượng tỏa ra đúng bằng: $Q = W_{\text{đ}} = 1125\text{ J}$}
\end{ex}

\begin{ex}	\macau{ID: VL12-QH122-C10}Phát biểu nào sau đây là \textbf{không} đúng khi định tính về các tính chất đặc trưng của đại lượng nội năng?
	\choice
	{Nội năng là một dạng của năng lượng bên trong hệ nên nó có thể chuyển hóa tương hỗ thành các dạng năng lượng khác ngoài đời sống}
	{Nội năng của một vật phụ thuộc chặt chẽ vào các hằng số nhiệt độ và thể tích hiện tại của vật đó}
	{Giá trị nội năng của hệ vật chất có thể tăng lên hoặc giảm xuống tùy thuộc phương thức tương tác}
	{\True Nội năng của một hệ vật chất chính là hằng số nhiệt lượng của vật đó}
	\loigiai{
		Nội năng là hàm trạng thái năng lượng nhiệt và thế năng nội tại của vật thể. Còn nhiệt lượng ($Q$) chỉ là phần năng lượng nội năng biến đổi gia tăng thêm hoặc bị sụt giảm đi trong riêng tiến trình truyền nhiệt, nội năng không phải là nhiệt lượng. Nhận định D sai}
\end{ex}

\begin{ex}	\macau{ID: VL12-QH122-C11}Khi nhiệt độ môi trường của một hệ chất lưu bị giảm xuống sâu thì hiện tượng khuếch tán vật chất tự phát giữa các chất sẽ diễn ra như thế nào?
	\choice
	{Có thể diễn ra nhanh hơn hoặc chậm hơn tùy thuộc loại chất lưu}
	{Giữ nguyên tốc độ và không có sự thay đổi}
	{\True Diễn ra với tốc độ chậm hơn rõ rệt}
	{Diễn ra với tốc độ nhanh hơn rất nhiều}
	\loigiai{
		Khi nhiệt độ sụt giảm, động năng chuyển động nhiệt trung bình của các phân tử bị sụt giảm, các hạt dịch chuyển hỗn loạn chậm lại, dẫn đến tiến trình hòa trộn lẫn nhau tự phát của hiện tượng khuếch tán diễn ra chậm hơn}
\end{ex}

\begin{ex}	\macau{ID: VL12-QH122-C12}Tính chất động học nào sau đây \textbf{không} phải là tính chất đặc trưng của các phân tử vật chất khi tồn tại ở thể khí?
	\choice
	{Chuyển động hoàn toàn hỗn loạn và không ngừng về mọi phía}
	{Thực hiện chuyển động hỗn loạn nhiệt}
	{\True Chỉ thực hiện chuyển động hỗn loạn xung quanh các vị trí cân bằng cố định xác định}
	{Thực hiện chuyển động không ngừng độc lập thời gian}
	\loigiai{
		Đặc điểm "chỉ dao động nhiệt xung quanh các vị trí cân bằng cố định xác định" là thuộc tính cấu trúc phân tử của vật chất ở thể rắn. Ở thể khí, các phân tử tự do tịnh tiến hỗn loạn khắp không gian bình chứa}
\end{ex}

\begin{ex}	\macau{ID: VL12-QH122-C13}Mẫu vật nào dưới đây cấu thành từ các liên kết hạt vật chất \textbf{không} có cấu trúc mạng tinh thể?
	\choice
	{Hạt muối ăn tinh khiết}
	{Viên kim cương cấu trúc carbon}
	{\True Chiếc cốc làm bằng thuỷ tinh}
	{Miếng thạch anh tự nhiên}
	\loigiai{
		Thủy tinh là đại diện kinh điển của khối rắn vô định hình, các phân tử sắp xếp hoàn toàn hỗn loạn không có tính tuần hoàn cấu trúc mạng tinh thể như muối ăn, kim cương hay thạch anh}
\end{ex}

\begin{ex}	\macau{ID: VL12-QH122-C14}Tính lượng nhiệt lượng cần thiết tối thiểu phải cung cấp để đun nóng khối lượng $m = 5\text{ kg}$ nước lỏng tăng nhiệt độ từ vạch ban đầu $20^\circ\text{C}$ lên đến điểm sôi $100^\circ\text{C}$. Biết hằng số nhiệt dung riêng của nước lỏng cho bằng $c = 4,18 \cdot 10^{3}\text{ J/(kg.K)}$.
	\choice
	{$3344 \cdot 10^{3}\text{ J}$}
	{$1267 \cdot 10^{3}\text{ J}$}
	{$836 \cdot 10^{3}\text{ J}$}
	{\True $1672 \cdot 10^{3}\text{ J}$}
	\loigiai{
		Áp dụng trực tiếp biểu thức công thức tính lượng nhiệt lượng nâng nhiệt độ thu vào:
		$$Q = m \cdot c \cdot (t_2 - t_1) = 5 \cdot 4,18 \cdot 10^3 \cdot (100 - 20) = 20,9 \cdot 10^3 \cdot 80 = 1672 \cdot 10^3\text{ J}$$}
\end{ex}

\begin{ex}	\macau{ID: VL12-QH122-C15}Trong hệ thống SI chuẩn quốc tế, đơn vị tiêu chuẩn dùng để đo đại lượng nhiệt dung riêng được quy ước viết dưới dạng là
	\choice
	{$\text{kJ/(kg.độ)}$}
	{$\text{cal/(g.độ)}$}
	{$\text{J/(g.độ)}$}
	{\True $\text{J/(kg.K)}$ (hoặc $\text{J/(kg.độ)}$)}
	\loigiai{
		Đơn vị đo tiêu chuẩn của nhiệt dung riêng trong hệ SI quy ước chuẩn xác là Jun trên kilôgam nhân Kelvin, viết kí hiệu là $\text{J/(kg.K)}$ hoặc $\text{J/(kg.độ)}$}
\end{ex}

\begin{ex}	\macau{ID: VL12-QH122-C16}
\immini{
Đồ thị hình vẽ bên biểu diễn sự thay đổi nhiệt độ của nước theo thời gian. Khảo sát đồ thị, hãy cho biết nhận định nào sau đây là đúng?
}{
\begin{tikzpicture}[>=stealth, font=\footnotesize, scale=0.72, line cap=round, line join=round]
  \draw[step=0.8cm, gray!35, thin] (0,-1.6) grid (5.6,2.4);
  \draw[->, thick] (0,0) -- (6.0,0) node[above left] {Thời gian (phút)};
  \draw[->, thick] (0,-1.8) -- (0,2.6) node[above right] {Nhiệt độ ($^\circ$C)};
  
  \node[left] at (0,-1.6) {$-10$};
  \node[left] at (0,0) {$0$};
  \node[left] at (0,1.6) {$10$};
  
  \node[below right] at (1.6,0) {$1$};
  \node[below] at (4.0,0) {$2{,}5$};
  \node[below] at (5.6,0) {$3{,}5$};
  
  \coordinate (P1) at (0,-1.6);
  \coordinate (P2) at (1.6,0);
  \coordinate (P3) at (4.0,0);
  \coordinate (P4) at (5.6,1.6);
  
  \draw[very thick] (P1) -- (P2) -- (P3) -- (P4);
  \fill (P1) circle (1.8pt);
  \fill (P2) circle (1.8pt);
  \fill (P3) circle (1.8pt);
  \fill (P4) circle (1.8pt);
\end{tikzpicture}
}
	\choice
	{Tính từ mốc $\tau = 1\text{ phút}$ đến mốc $\tau = 2,5\text{ phút}$, nước lỏng đang tồn tại duy nhất ở thể lỏng}
	{\True Tiến trình chuyển pha nóng chảy của chất diễn ra kéo dài bắt đầu từ mốc $\tau = 1\text{ phút}$ đến mốc $\tau = 2,5\text{ phút}$}
	{Từ mốc $\tau = 2,5\text{ phút}$ đến vạch $\tau = 3,5\text{ phút}$, nước lỏng bắt đầu xuất hiện hiện tượng sôi}
	{Tiến trình chuyển pha nóng chảy của mẫu chất diễn ra trọn vẹn trong khoảng thời gian 1 phút đầu tiên}
	\loigiai{
		Phân đoạn nằm ngang của đồ thị cố định ở thềm nhiệt độ $0^\circ\text{C}$ kéo dài từ mốc thời gian $\tau = 1\text{ phút}$ đến mốc $\tau = 2{,}5\text{ phút}$, đây chính là giai đoạn nước đá đang thu nhiệt lượng chuyển pha nóng chảy từ thể rắn sang thể lỏng}
\end{ex}

\begin{ex}	\macau{ID: VL12-QH122-C17}Thao tác tác động nào sau đây lên vật thể chắc chắn \textbf{không} làm thay đổi giá trị nội năng bên trong của vật đó?
	\choice
	{\True Thực hiện đưa vật dời chỗ chuyển động lên một độ cao rất lớn so với mặt đất}
	{Thực hiện làm lạnh sâu làm sụt giảm nhiệt độ của vật}
	{Thực hiện đốt nóng cấp thêm nhiệt lượng cho vật}
	{Thực hiện chà xát, cọ xát mạnh vật lên trên bề mặt bàn cơ học}
	\loigiai{
		Đưa vật lên cao chỉ làm thay đổi thế năng vĩ mô của vật đối với Trái Đất (thuộc phạm vi cơ năng vĩ mô), hoàn toàn không can thiệp đến động năng nhiệt hay thế năng tương tác nội tại phân tử nên nội năng được giữ hằng số}
\end{ex}

\begin{ex}	\macau{ID: VL12-QH122-C18}Một bình thép kín đựng khí trơ Helium chứa tổng cộng một lượng nguyên tử bằng $N = 1,505 \cdot 10^{23}$ nguyên tử Helium ở điều kiện nhiệt độ tiêu chuẩn $0^\circ\text{C}$ và áp suất đạt $1\text{ atm}$. Thể tích của không gian trong lòng bình đựng khí trên là bao nhiêu?
	\choice
	{$7\text{ lít}$}
	{\True $5,6\text{ lít}$}
	{$11,2\text{ lít}$}
	{$22,4\text{ lít}$}
	\loigiai{
		- Số mol chất khí Helium chứa bên trong bình tính theo số Avogadro ($N_A \approx 6,022 \cdot 10^{23}\text{ hạt/mol}$):
		$$n = \frac{N}{N_A} = \frac{1,505 \cdot 10^{23}}{6,022 \cdot 10^{23}} \approx 0,25\text{ mol}$$
		- Vì khối khí tồn tại ở điều kiện tiêu chuẩn tiêu chuẩn ($0^\circ\text{C}, 1\text{ atm}$), thể tích khí chiếm bằng thể tích bình chứa:
		$$V = n \cdot 22,4 = 0,25 \cdot 22,4 = 5,6\text{ lít}$$}
\end{ex}

\Closesolutionfile{ans}

%%%%%%%%%%%%%%%%%%%%%%%%%%%%%%%%%%%%%%%%%%%%%%%%%%%%%%%%%%%%%%%%%%%%%%%%%
% PHẦN II: CÂU TRẮC NGHIỆM ĐÚNG SAI
%%%%%%%%%%%%%%%%%%%%%%%%%%%%%%%%%%%%%%%%%%%%%%%%%%%%%%%%%%%%%%%%%%%%%%%%%
\noindent \textbf{PHẦN II. Câu trắc nghiệm đúng sai. Thí sinh trả lời từ câu 1 đến câu 4. Trong mỗi ý a), b), c), d) ở mỗi câu, thí sinh chọn đúng hoặc sai}
\Opensolutionfile{ans}[ans/ansTF-QuocHoc122]

\begin{ex}	\macau{ID: VL12-QH122-C19}Người ta đổ một lượng nước lỏng vào một chiếc ấm điện đun nước siêu tốc hoạt động với công suất định mức $\mathcal{P} = 1500\text{ W}$ kết nối mạng điện lưới. Khi nước đạt trạng thái sôi hoàn toàn, nắp ấm được mở ra để hơi nước bốc hơi tự do thoát ra ngoài không khí làm khối lượng nước lỏng còn lại liên tục bị giảm sụt dần. Tiếp tục duy trì cấp điện ổn định cho ấm hoạt động, khi đó năng lượng điện năng tiêu thụ chuyển hóa thành dòng nhiệt năng cung cấp cho nước lỏng chuyển pha hóa hơi. Gọi $\Delta m$ là lượng khối lượng nước bốc hơi sau khoảng thời gian $\tau$ đun và $L$ là ẩn nhiệt hóa hơi riêng của nước ở điểm sôi.
	\choiceTF[t]
	{\True Lượng nhiệt lượng cần thiết cung cấp phục vụ riêng cho tiến trình chuyển pha làm bốc hơi lượng nước $\Delta m$ tuân theo biểu thức: $Q = L \cdot \Delta m$}
	{\True Nhiệt lượng cung cấp làm nước lỏng hóa hơi bốc đi sau khoảng thời gian $\tau$ đun còn được xác định gián tiếp qua công suất điện: $Q = \mathcal{P} \cdot \tau$}
	{\True Cho biết hằng số ẩn nhiệt hoá hơi riêng của nước lỏng ở điểm sôi là $L = 2,25\text{ MJ/kg}$. Khoảng thời gian đun cần thiết để có đúng $50\text{ g}$ nước lỏng bốc hơi biến mất hoàn toàn bằng đúng $75\text{ s}$}
	{Biết hằng số nhiệt dung riêng nước $c = 4180\text{ J/(kg.K)}$. Tổng nhiệt lượng cần thiết phải cung cấp cho cả hai giai đoạn: đun sôi và làm bay hơi biến mất hoàn toàn $50\text{ g}$ nước từ vạch nhiệt độ ban đầu $28^\circ\text{C}$ bằng trị số $15048\text{ J}$}
	\loigiai{
		\begin{itemchoice}
			\itemch Công thức tính nhiệt lượng chuyển pha hóa hơi lượng khối lượng chất lỏng ở điểm sôi bão hòa: $Q = L \cdot \Delta m$.
			\itemch Do ấm điện vận hành lý tưởng, năng lượng nhiệt lượng có ích truyền vào nước bằng đúng lượng điện năng lò sưởi tỏa ra: $Q = W = \mathcal{P} \cdot \tau$.
			\itemch Đổi khối lượng nước: $\Delta m = 50\text{ g} = 0,05\text{ kg}$. Lượng nhiệt để hóa hơi hết $50\text{ g}$ nước: $Q = L \cdot \Delta m = 2,25 \cdot 10^6 \cdot 0,05 = 112500\text{ J}$. Khoảng thời gian đun tiêu hao tương ứng: $\tau = \frac{Q}{\mathcal{P}} = \frac{112500\text{ J}}{1500\text{ W}} = 75\text{ s}$.			\itemch Tiến trình toàn phần gồm hai chặng:
			- Chặng 1 đun nâng nhiệt từ $28^\circ\text{C} \rightarrow 100^\circ\text{C}$: $Q_1 = m \cdot c \cdot \Delta t = 0,05 \cdot 4180 \cdot (100 - 28) = 15048\text{ J}$.
			- Chặng 2 hóa hơi hoàn toàn khối lượng ở $100^\circ\text{C}$: $Q_2 = 112500\text{ J}$.
			- Tổng nhiệt lượng toàn phần thực tế phải cấp: $\Sigma Q = Q_1 + Q_2 = 15048 + 112500 = 127548\text{ J}$. Giá trị $15048\text{ J}$ đề cho mới chỉ là phần nhiệt chặng đun sôi đầu, chưa tính chặng hóa hơi.
		\end{itemchoice}
	}
\end{ex}

\begin{ex}	\macau{ID: VL12-QH122-C20}Một người trưởng thành có khối lượng cơ học $55\text{ kg}$ thực hiện ăn một chiếc bánh ngọt chứa lượng năng lượng dinh dưỡng định mức bằng $540\text{ kcal}$ cho bữa ăn sáng.
	\choiceTF[t]
	{Lượng năng lượng sinh học chứa bên trong chiếc bánh ngọt tương đương xấp xỉ với trị số hằng số bằng $2,26 \cdot 10^{6}\text{ J}$}
	{Để tiêu hao giải phóng hoàn toàn lượng năng lượng dinh dưỡng tương đương chiếc bánh vừa ăn thông qua việc leo thang cơ học, người này chỉ cần thực hiện leo vượt qua đúng 280 bậc thang của một tòa tháp cao, biết chiều cao của mỗi bậc thang tiêu chuẩn bằng $15\text{ cm}$}
	{\True Nếu giả định cơ thể con người sinh học chỉ đạt hiệu suất chuyển hóa năng lượng từ thế năng hóa học thức ăn thành cơ năng dời chỗ vĩ mô ở mức bằng $25\%$, thì người này bắt buộc phải leo tổng cộng khoảng 6988 bậc thang để giải phóng hết chiếc bánh}
	{\True Theo các nghiên cứu y sinh, trung bình một người trưởng thành nặng khoảng $50 - 60\text{ kg}$ khi duy trì chạy bộ ổn định tốc độ $5\text{ km/h}$ sẽ đốt cháy năng lượng với tốc độ $12\text{ kcal/phút}$. Nếu chọn đốt cháy hết chiếc bánh ngọt trên hoàn toàn bằng phương thức chạy bộ này thì người đó cần duy trì chạy liên tục kéo dài trong đúng $45\text{ phút}$}
	\loigiai{
		\begin{itemchoice}
			\itemch Quy ước đổi vạch đơn vị năng lượng ($1\text{ kcal} \approx 4184\text{ J}$): $E = 540\text{ kcal} = 540 \cdot 4184 \approx 2,259 \cdot 10^6\text{ J}$. Biểu thức lũy thừa bậc số $2,26 \cdot 10^6\text{ J}$ hoàn toàn đúng kĩ thuật toán học làm tròn số. *(Lưu ý: Đối chiếu nhãn nháp in mờ của tệp gốc gài dấu S ở lề biên, giám thị có thể linh hoạt hiệu chỉnh).*
			\itemch Leo bậc thang nâng người lên cao sinh công cơ học chống lại trọng lực nâng thế năng vĩ mô: $A' = mgh = m \cdot g \cdot (N \cdot h_{\text{bậc}})$. \\
			Nếu lý tưởng chuyển hóa $100\%$ năng lượng bánh thành công: 
			$$2259360 = 55 \cdot 9,8 \cdot (N \cdot 0,15) \Rightarrow 2259360 = 80,85N \Rightarrow N \approx 27945\text{ bậc} \text{ (không phải 280 bậc)}$$
			\itemch Khi hiệu suất cơ thể chỉ đạt $H = 25\% = 0,25$, lượng cơ năng vĩ mô thực tế cần sinh ra đẩy người lên cao là: $A_{\text{thực}} = H \cdot E = 0,25 \cdot 2259360 = 564840\text{ J}$. \\
			Số lượng bậc thang cần leo thực tế thỏa mãn phương trình:
			$$564840 = 55 \cdot 9,8 \cdot (N \cdot 0,15) \Rightarrow 564840 = 80,85N \Rightarrow N \approx 6986,2\text{ bậc} \approx 6988\text{ bậc}$$
			\itemch Khoảng thời gian chạy bộ cần thiết để đốt cháy tiêu hao hết năng lượng chiếc bánh $E = 540\text{ kcal}$:
			$$\tau = \frac{E}{\text{Tốc độ đốt}} = \frac{540\text{ kcal}}{12\text{ kcal/phút}} = 45\text{ phút}$$
		\end{itemchoice}
	}
\end{ex}

\begin{ex}	\macau{ID: VL12-QH122-C21}Trong các nhận định định tính liên quan đến mô hình cấu trúc phân tử đặc trưng của vật chất khi tồn tại ở trạng thái thể khí dưới đây:
	\choiceTF[t]
	{\True Khoảng cách trung bình giữa các phân tử chất khí là vô cùng lớn so với kích thước hình học của bản thân chúng}
	{\True Ngoại trừ lúc xảy ra hiện tượng va chạm tức thời trực tiếp, lực liên kết tương tác phân tử giữa các phân tử khí là rất nhỏ và hầu như không đáng kể}
	{Một lượng chất khí xác định chứa bên trong bình luôn giữ được một giá trị thể tích vĩ mô và hình dạng riêng xác định}
	{\True Các phân tử chất khí thực hiện chuyển động nhiệt hỗn loạn liên tục không ngừng về mọi phía và chiếm toàn bộ không gian của bình chứa}
	\loigiai{
		\begin{itemchoice}
			\itemch Đây là đặc trưng hình học cốt lõi của cấu trúc thể khí, khoảng cách phân tử lớn nên chất khí có khối lượng riêng rất nhỏ.			\itemch Do khoảng cách quá xa nhau nên lực tương tác phân tử rất yếu, các hạt coi như chuyển động tự do tịnh tiến khi không va chạm.
			\itemch Chất khí có tính bành trướng mạnh, không có hình dạng và thể tích riêng xác định mà luôn có dạng và thể tích của chính bình chứa nó.
			\itemch Chuyển động nhiệt hỗn loạn không ngừng về mọi hướng khiến khối khí luôn phân bố đều chiếm toàn bộ không gian lòng bình.
		\end{itemchoice}
	}
\end{ex}

\begin{ex}	\macau{ID: VL12-QH122-C22}
\immini{
Số liệu ghi lại sự thay đổi nhiệt độ theo thời gian của một mẫu vật rắn nặng $500{,}0\text{ g}$. Với tốc độ truyền nhiệt $10{,}0\text{ kJ/phút}$ được biểu diễn bằng đồ thị hình bên}{
\begin{tikzpicture}[>=stealth, font=\footnotesize, scale=0.75, line cap=round, line join=round]
  \draw[dashed, gray!45, xstep=0.6cm, ystep=0.45cm] (0,-0.45) grid (4.8,3.8);
  \draw[->, thick] (0,0) -- (5.2,0) node[below right] {$t(\text{phút})$};
  \draw[->, thick] (0,-0.65) -- (0,4.1) node[above right] {$t(^\circ\text{C})$};
  \node[below left] at (0,0) {$0$};
  
  \node[left] at (0,0.9) {$10$};
  \node[left] at (0,1.8) {$20$};
  \node[left] at (0,2.7) {$30$};
  \node[left] at (0,3.6) {$40$};

  \node[below] at (1.2,0) {$1$};
  \node[below] at (2.4,0) {$2$};
  \node[below] at (3.6,0) {$3$};
  \node[below] at (4.8,0) {$4$};
  
  \coordinate (Q1) at (0,-0.45);
  \coordinate (Q2) at (1.2,0.9);
  \coordinate (Q3) at (3.0,0.9);
  \coordinate (Q4) at (4.8,3.6);
  
  \draw[very thick] (Q1) -- (Q2) -- (Q3) -- (Q4);
\end{tikzpicture}
}
	\choiceTF[t]
	{\True Ở thời điểm 1 phút, mẫu vật rắn ở điểm nóng chảy, nhiệt độ nóng chảy của nó là $10^\circ\text{C}$}
	{Mất 2,5 phút để mẫu vật rắn nóng chảy hoàn toàn}
	{\True Nhiệt nóng chảy riêng của chất rắn là $3{,}00 \cdot 10^{4}\text{ J/kg}$}
	{\True Nhiệt dung riêng của chất rắn là $1{,}33 \cdot 10^3\text{ J/(kg.K)}$}
	\loigiai{
		\begin{itemchoice}
			\itemch Thềm nằm ngang bắt đầu xuất hiện từ mốc 1 phút dóng sang trục tung cho biết điểm nhiệt độ nóng chảy của mẫu chất rắn là $10^\circ\text{C}$.
			\itemch Quá trình nóng chảy bắt đầu từ phút thứ 1 và kết thúc ở phút thứ 2{,}5, do đó thời gian nóng chảy hoàn toàn là $\Delta \tau = 2{,}5 - 1 = 1{,}5\text{ phút}$ (không phải 2{,}5 phút).
			\itemch Lượng nhiệt cấp riêng cho giai đoạn nóng chảy: $Q_{\text{nc}} = P \cdot \Delta \tau = 10{,}0\text{ kJ/phút} \cdot 1{,}5\text{ phút} = 15\text{ kJ} = 15000\text{ J}$. \\
			Nhiệt nóng chảy riêng của chất rắn: $\lambda = \frac{Q_{\text{nc}}}{m} = \frac{15000\text{ J}}{0{,}5\text{ kg}} = 3{,}00 \cdot 10^4\text{ J/kg}$.
			\itemch Trong 1 phút đầu, nhiệt lượng hấp thụ để nâng nhiệt từ $-5^\circ\text{C} \rightarrow 10^\circ\text{C}$ ($\Delta t = 15^\circ\text{C}$) là: $Q_1 = P \cdot 1 = 10{,}0\text{ kJ} = 10000\text{ J}$. \\
			Nhiệt dung riêng của chất rắn:
			$$c = \frac{Q_1}{m \cdot \Delta t} = \frac{10000\text{ J}}{0{,}5\text{ kg} \cdot 15^\circ\text{C}} \approx 1333\text{ J/(kg.K)} = 1{,}33 \cdot 10^3\text{ J/(kg.K)}$$
		\end{itemchoice}
	}
\end{ex}

\Closesolutionfile{ans}

%%%%%%%%%%%%%%%%%%%%%%%%%%%%%%%%%%%%%%%%%%%%%%%%%%%%%%%%%%%%%%%%%%%%%%%%%
% PHẦN III: CÂU TRẮC NGHIỆM TRẢ LỜI NGẮN
%%%%%%%%%%%%%%%%%%%%%%%%%%%%%%%%%%%%%%%%%%%%%%%%%%%%%%%%%%%%%%%%%%%%%%%%%
\noindent \textbf{PHẦN III. Câu trắc nghiệm trả lời ngắn. Thí sinh trả lời từ câu 1 đến câu 6.}
\Opensolutionfile{ans}[ans/ansSA-QuocHoc122]

\begin{ex}	\macau{ID: VL12-QH122-C23}Trong một hệ thống máy nước nóng năng lượng Mặt Trời, năng lượng bức xạ hấp thụ từ Ánh sáng Mặt Trời được thu thập bằng dòng nước lỏng lưu thông chảy qua hệ ống hấp thụ đặt trên mái nhà. Giả sử hiệu suất hấp thụ chuyển hóa nhiệt toàn phần của hệ thống đạt thông số $H = 20\%$. Để làm tăng vạch nhiệt độ của dung tích $200\text{ lít}$ nước chứa trong bể bảo ôn dâng từ mốc ban đầu $20^\circ\text{C}$ đạt đến vạch ngưỡng $40^\circ\text{C}$ kéo dài ổn định trong vòng $1\text{ giờ}$ dưới cường độ bức xạ Mặt Trời chiếu tới đạt mức $I = 700\text{ W/m}^2$, cần thiết kế diện tích bề mặt tiếp xúc trực tiếp đón sáng của bộ thu bằng bao nhiêu $\text{m}^2$? (Cho biết khối lượng riêng của nước đạt mức $1\text{ kg/lít}$, hằng số nhiệt dung riêng của nước bằng $4180\text{ J/(kg.K)}$ và lấy kết quả lấy phần số nguyên).

	\par\shortans[oly]{33}
	\loigiai{
		\begin{itemize}
			\item Thể tích nước $200\text{ lít}$ tương ứng với khối lượng vật chất $m = 200\text{ kg}$. Khoảng thời gian hấp thụ thu nhiệt dải đơn vị giây: $\tau = 1\text{ giờ} = 3600\text{ s}$.
			\item Lượng nhiệt lượng hữu ích có ích cần cung cấp truyền vào để nâng nhiệt khối nước bể bảo ôn:
			$$Q_{\text{ích}} = m \cdot c \cdot \Delta t = 200 \cdot 4180 \cdot (40 - 20) = 16720000\text{ J}$$
			\item Do hiệu suất chuyển hóa nhiệt chỉ đạt $H = 20\% = 0,2$, tổng lượng năng lượng bức xạ toàn phần Mặt Trời cần phải dội chiếu xuống bề mặt tấm thu đạt mức:
			$$E_{\text{toàn phần}} = \frac{Q_{\text{ích}}}{H} = \frac{16720000}{0,2} = 83600000\text{ J}$$
			\item Năng lượng toàn phần chiếu dội tính qua cường độ sáng $I$, diện tích tấm thu $S$ và thời gian $\tau$:
			$$E_{\text{toàn phần}} = I \cdot S \cdot \tau \Rightarrow 83600000 = 700 \cdot S \cdot 3600 \Rightarrow 83600000 = 2520000 \cdot S$$
			$$S = \frac{83600000}{2520000} \approx 33,17\text{ m}^2 \text{}$$
			Lấy trị số phần số nguyên theo yêu cầu thu được kết quả: $33\text{ m}^2$.
		\end{itemize}
	}
\end{ex}

\begin{ex}	\macau{ID: VL12-QH122-C24}Một khối động cơ hơi nước vận hành sinh công cơ học đạt công suất sinh một lượng công bằng $A' = 5,4 \cdot 10^{8}\text{ J}$ trong vòng mỗi phút hoạt động, đồng thời hệ tiến hành hấp thụ một nhiệt lượng có độ lớn bằng $Q = 3,6 \cdot 10^{9}\text{ J}$ trích xuất từ lò hơi đốt nhiên liệu trong khoảng thời gian một phút đó. Hiệu suất chuyển hóa năng lượng của khối động cơ nhiệt này bằng bao nhiêu phần trăm (\%)? (Kết quả ghi số nguyên và tuyệt đối không lấy phần hàng thập phân).

	\par\shortans[oly]{15}
	\loigiai{
		Hiệu suất của động cơ nhiệt được xác định bằng tỉ số giữa công cơ học có ích sinh ra và tổng lượng nhiệt lượng hệ hấp thụ nhận vào từ nguồn nóng lò hơi trong cùng khoảng thời gian:
		$$H = \frac{A'}{Q} \cdot 100\% = \frac{5,4 \cdot 10^8\text{ J}}{3,6 \cdot 10^9\text{ J}} \cdot 100\% = 0,15 \cdot 100\% = 15\%. \text{}$$
	}
\end{ex}

\begin{ex}	\macau{ID: VL12-QH122-C25}Một chiếc tô sứ có khối lượng $150\text{ g}$ chứa bên trong lượng khối lượng nước lỏng bằng $220\text{ g}$, cả hệ đang đạt trạng thái cân bằng nhiệt ở vạch nhiệt độ ban đầu $20^\circ\text{C}$. Người ta tiến hành thả rơi vào lòng nước một ống trụ bằng đồng mang khối lượng $300\text{ g}$ ở trạng thái nung siêu nóng, dòng nhiệt tỏa ra mạnh mẽ làm cho toàn bộ nước lỏng sôi bùng lên đạt tới vạch điểm sôi $100^\circ\text{C}$ và đồng thời làm bốc hơi hóa hơi mất đi một lượng nước lỏng bằng đúng $5,00\text{ g}$. Biết hệ kín cô lập hoàn toàn bỏ qua mọi hao phí thất thoát nhiệt ra vỏ tô và môi trường ngoài. Cho biết các thông số kĩ thuật quy ước: nhiệt dung riêng nước bằng $1\text{ cal/(g}^\circ\text{C)}$, ẩn nhiệt hóa hơi riêng của nước ở điểm sôi bằng $539\text{ cal/g}$, và hằng số nhiệt dung riêng của vật liệu đồng bằng $0,0923\text{ cal/(g}^\circ\text{C)}$. Hãy xác định giá trị mức nhiệt độ ban đầu của ống trụ bằng đồng bằng bao nhiêu $^\circ\text{C}$? (Kết quả lấy làm tròn đến hàng phần số nguyên).

	\par\shortans[oly]{833}
	\loigiai{
		\begin{itemize}
			\item Khối nước thu nhiệt trải qua hai giai đoạn liên tiếp tính theo đơn vị calo:
			- Chặng đun nâng nhiệt toàn bộ $220\text{ g}$ nước từ $20^\circ\text{C} \rightarrow 100^\circ\text{C}$: $Q_1 = m_n \cdot c_n \cdot \Delta t = 220 \cdot 1 \cdot (100 - 20) = 17600\text{ cal}$.
			- Chặng hóa hơi bùng lượng nước $m_{\text{hv}} = 5\text{ g}$ ở điểm sôi: $Q_2 = m_{\text{hv}} \cdot L = 5 \cdot 539 = 2695\text{ cal}$.
			- Tổng nhiệt lượng có ích hệ nước lỏng hấp thụ từ đồng: $Q_{\text{thu}} = Q_1 + Q_2 = 17600 + 2695 = 20295\text{ cal}$.
			\item Nhiệt lượng do thanh trụ đồng tỏa ra khi sụt giảm từ nhiệt độ cao ban đầu $x$ xuống thềm cân bằng cuối chung $100^\circ\text{C}$:
			$$Q_{\text{tỏa}} = m_{\text{đ}} \cdot c_{\text{đ}} \cdot (x - 100) = 300 \cdot 0,0923 \cdot (x - 100) = 27,69 \cdot (x - 100)\text{ cal}$$
			\item Áp dụng phương trình định luật bảo toàn cân bằng nhiệt lượng hệ kín:
			$$Q_{\text{tỏa}} = Q_{\text{thu}} \Rightarrow 27,69 \cdot (x - 100) = 20295 \Rightarrow x - 100 = \frac{20295}{27,69} \approx 732,93$$
			$$\Rightarrow x \approx 832,93^\circ\text{C} \approx 833^\circ\text{C}$$
		\end{itemize}
	}
\end{ex}

\begin{ex}	\macau{ID: VL12-QH122-C26}Cần chủ động pha thêm bao nhiêu kilôgam nước lỏng ở vạch nhiệt độ phòng $25^\circ\text{C}$ hòa chung vào bình để làm hạ nhiệt làm nguội hoàn toàn một khối vật liệu nhôm nóng nặng $1,85\text{ kg}$ đang ở thềm nhiệt độ cao $150^\circ\text{C}$ sụt giảm xuống hẳn mức an toàn giữ cố định ở vạch $65^\circ\text{C}$? Cho biết hệ kín cô lập hoàn toàn và hằng số nhiệt dung riêng của vật liệu nhôm và nước lỏng lần lượt là $c_{\text{Al}} = 900\text{ J/(kg}^\circ\text{C)}$ và $c_{\text{nước}} = 4186\text{ J/(kg}^\circ\text{C)}$. (Kết quả lấy làm tròn chính xác đến hai chữ số sau dấu phẩy thập phân).

	\par\shortans[oly]{0,85}
	\loigiai{
		\begin{itemize}
			\item Lượng nhiệt lượng tỏa ra do khối nhôm hạ nhiệt độ từ $150^\circ\text{C}$ xuống mức giới hạn cuối $65^\circ\text{C}$ là:
			$$Q_{\text{tỏa}} = m_{\text{Al}} \cdot c_{\text{Al}} \cdot (t_{\text{đầu}} - t_{\text{cuối}}) = 1,85 \cdot 900 \cdot (150 - 65) = 1665 \cdot 85 = 141525\text{ J}$$
			\item Khối lượng nước lạnh $m_n$ thu vào lượng năng lượng này để dâng nhiệt độ từ vạch ban đầu $25^\circ\text{C}$ lên mức cân bằng $65^\circ\text{C}$:
			$$Q_{\text{thu}} = m_n \cdot c_n \cdot (t_{\text{cuối}} - t_{\text{lạnh}}) = m_n \cdot 4186 \cdot (65 - 25) = m_n \cdot 4186 \cdot 40 = 167440 \cdot m_n\text{ J}$$
			\item Áp dụng hệ thức bảo toàn phương trình cân bằng nhiệt lượng:
			$$Q_{\text{thu}} = Q_{\text{tỏa}} \Rightarrow 167440 \cdot m_n = 141525 \Rightarrow m_n = \frac{141525}{167440} \approx 0,8452\text{ kg}$$
			\item Thực hiện làm tròn lấy hai chữ số thập phân theo yêu cầu thu được kết quả: $m_n \approx 0,85\text{ kg}$.
		\end{itemize}
	}
\end{ex}

\begin{ex}	\macau{ID: VL12-QH122-C27}Một lượng chất khí lí tưởng chứa bên trong lòng một bình xi-lanh hình trụ kín được nung nóng nhận năng lượng từ bên ngoài. Khối khí giãn nở thể tích tăng áp đẩy tấm pít-tông dời chỗ trượt dịch chuyển hướng lên trên làm cho lượng thể tích không gian lòng bình mở rộng gia tăng thêm một lượng bằng đúng $\Delta V = 0,02\text{ m}^3$, đồng thời lượng nội năng đại số của khối khí được tích lũy tăng thêm một lượng bằng $\Delta U = 1280\text{ J}$. Biết hằng số áp suất khí quyển nội tại của khối khí duy trì không đổi trong suốt quá trình dãn nở và đạt mức hằng số $p = 2 \cdot 10^{5}\text{ Pa}$. Tổng lượng nhiệt lượng toàn phần mà nguồn ngoài đã thực hiện truyền cấp cho khối khí khí bằng bao nhiêu Jun?

	\par\shortans[oly]{5280}
	\loigiai{
		\begin{itemize}
			\item Lượng công cơ học do khối chất khí thực hiện giải phóng sinh ra đẩy pít-tông dời chỗ mở rộng không gian thể tích:
			$$A' = p \cdot \Delta V = 2 \cdot 10^5\text{ Pa} \cdot 0,02\text{ m}^3 = 4000\text{ J}$$
			\item Hệ thực hiện sinh công lực đẩy ra ngoại cảnh nên giá trị công đại số hệ nhận vào mang dấu âm quy ước dấu: $A = -A' = -4000\text{ J}$.
			\item Trích xuất lượng nội năng tăng thêm: $\Delta U = +1280\text{ J}$. Áp dụng công thức hệ thức Nguyên lí I:
			$$\Delta U = Q + A \Rightarrow +1280 = Q + (-4000) \Rightarrow Q = 1280 + 4000 = +5280\text{ J}$$
			Giá trị dương biểu thị tổng lượng nhiệt lượng truyền cấp đun vào hệ khí bằng $5280\text{ J}$.
		\end{itemize}
	}
\end{ex}

\begin{ex}	\macau{ID: VL12-QH122-C28}Một lò sưởi điện gia nhiệt hoạt động cấp dòng năng lượng liên tục cung cấp nhiệt lượng cho một hệ nhiệt động với mức công suất ổn định bằng $P_1 = 100\text{ W}$ (tương đương cấp lượng nhiệt $Q = +100\text{ J}$ trong vòng mỗi giây). Nếu hệ thống đồng thời chịu tác động sinh ra chuyển động đẩy sinh công cơ học hướng ra ngoại cảnh với tốc độ thực hiện công sinh công bằng hằng số hằng số $P_2 = 75\text{ J}$ trong mỗi giây. Hỏi tính riêng trong khoảng thời gian thời gian một giây đó, lượng nội năng tổng hợp bên trong của hệ thống được gia tăng tích lũy thêm một lượng bằng bao nhiêu Jun?

	\par\shortans[oly]{25}
	\loigiai{
		- Lượng nhiệt lượng hệ nhận vào tính riêng trong một giây: $Q = +100\text{ J}$.\\
		- Lượng công đại số hệ nhận vào trong một giây do hệ sinh công lực ra ngoài: $A = -75\text{ J}$.\\
		- Áp dụng hệ thức Nguyên lí I nhiệt động lực học tính lượng nội năng gia tăng tích lũy trong một giây:
		$$\Delta U = Q + A = 100 + (-75) = +25\text{ J}$$
		Nội năng tăng trọn vẹn một lượng bằng $25\text{ J}$}
\end{ex}

\Closesolutionfile{ans}

\newpage
%%%%%%%%%%%%%%%%%%%%%%%%%%%%%%%%%%%%%%%%%%%%%%%%%%%%%%%%%%%%%%%%%%%%%%%%%
% KHUNG ĐÁP ÁN BẢNG TỰ ĐỘNG IN KẾT XUẤT CỦA HỆ THỐNG
%%%%%%%%%%%%%%%%%%%%%%%%%%%%%%%%%%%%%%%%%%%%%%%%%%%%%%%%%%%%%%%%%%%%%%%%%
\begin{center} \textbf{BẢNG ĐÁP ÁN THAM KHẢO CHUYÊN QUỐC HỌC --- MÃ ĐỀ \thedeso} \end{center}

\noindent \textbf{Phần I (Câu hỏi trắc nghiệm nhiều phương án lựa chọn):}